
\documentclass[11pt]{article}

% ── Packages ──────────────────────────────────────────────────────────────────
\usepackage[margin=1in]{geometry}
\usepackage{amsmath, amssymb, amsthm}
\usepackage{mathtools}
\usepackage{booktabs}
\usepackage{hyperref}
\usepackage{cleveref}
\usepackage{xcolor}
\usepackage{algorithm}
\usepackage{algpseudocode}
\usepackage{graphicx}
\usepackage{subcaption}
\graphicspath{{../experiments/demand_understanding/}}

% ── Notation shorthands ───────────────────────────────────────────────────────
\newcommand{\R}{\mathbb{R}}
\newcommand{\E}{\mathbb{E}}
\newcommand{\cU}{\mathcal{U}}
\newcommand{\cS}{\mathcal{S}}
\DeclareMathOperator*{\argmin}{arg\,min}

\theoremstyle{definition}
\newtheorem{definition}{Definition}
\newtheorem{remark}{Remark}
\newtheorem{proposition}{Proposition}

% ── Title ─────────────────────────────────────────────────────────────────────
\title{%
  Robust Station Selection for Microtransit\\
  Under Demand Uncertainty%
}
\author{Gabriel Yong}
\date{}

% ──────────────────────────────────────────────────────────────────────────────
\begin{document}
\maketitle

\section{Introduction}

Microtransit is a form of transit that seeks to fill the gap between existing services such as private-hire vehicles and scheduled public transit. It is an on-demand, publicly accessible shared-ride service. Like private-hire, it is flexible and on-demand, but it also benefits from the economies of scale of scheduled public transportation by using mid-size capacity vehicles and allowing for greater vehicle utilization.

However, having a fully flexible microtransit system is overly complex and reduces its aggregation capabilities, often having to make significant detours just to satisfy one or two customers. Thus, selection of these \emph{virtual bus stops} (VBS) becomes an important strategic decision. The optimal selection enables better passenger aggregation, reducing total vehicle distance travelled and improving vehicle utilization, while offering passengers decreased waiting times and fewer detours than scheduled public transit.

While the VBS selection is able to adapt between service periods within the day (morning, afternoon, evening, night), the selection for each service period should be robust against shifts in the demand patterns. Within a service period, demand can fluctuate from day to day and the system should still provide a relatively similar experience across each day. We thus address this problem with a two-stage robust model where origin-destination (OD) demand varies within a budgeted uncertainty set.

\section{Model}

We formulate VBS selection as a two-stage mixed-integer program. The first stage designates a fixed set of $l$ stations shared across all service periods. The second stage activates $k$ of these per period, and OD pair assignments are used to evaluate the demand-weighted walking and routing cost of each activation plan. We first present the nominal model assuming known demand, then characterise the demand uncertainty set, and finally derive the robust counterpart.

\subsection{Nominal Model}

\paragraph{Sets.}
\[
	\begin{aligned}
		J      &  &  & \text{set of candidate stations,}                                  \\
		S      &  &  & \text{set of service scenarios,}                                   \\
		\Omega &  &  & \text{set of origin-destination pairs,}                            \\
		A_{od} &  &  & \text{admissible pickup--dropoff station pairs for OD pair }(o,d),
	\end{aligned}
\]
where
\[
	A_{od}
	=
	\left\{
	(j,k)\in J\times J :
	d^{\mathrm{orig}}_{oj}\le W^{\max},\
	d^{\mathrm{dest}}_{dk}\le W^{\max}
	\right\}.
\]
In the experimental instance, $A_{od}$ is non-empty for all $(o,d)\in\Omega$.

\paragraph{Input parameters.}
The following parameters are assumed to be known:
\[
	\begin{aligned}
		q_{ods}                & \ge 0              &  & \text{demand for OD pair }(o,d)\text{ in scenario }s,              \\
		d^{\mathrm{orig}}_{oj} & \ge 0              &  & \text{walking distance from origin }o\text{ to station }j,         \\
		d^{\mathrm{dest}}_{dk} & \ge 0              &  & \text{walking distance from station }k\text{ to destination }d,    \\
		c_{jk}                 & \ge 0              &  & \text{vehicle travel cost from station }j\text{ to station }k,     \\
		\lambda                & \ge 0              &  & \text{weight on vehicle travel cost relative to walking distance,} \\
		W^{\max}               & \ge 0              &  & \text{maximum walking distance allowed for a passenger,}           \\
		l                      & \in \mathbb{Z}_{+} &  & \text{number of stations to select,}                               \\
		k                      & \in \mathbb{Z}_{+} &  & \text{number of stations to activate per scenario.}
	\end{aligned}
\]

\paragraph{Decision variables.}
The model has three layers of binary decisions:
\begin{align*}
	y_j       & = \begin{cases} 1 & \text{station } j \text{ is selected } \\ 0 & \text{otherwise} \end{cases}                                                                               \\[6pt]
	z_{js}    & = \begin{cases} 1 & \text{station } j \text{ is activated and available for service in scenario } s \\ 0 & \text{otherwise} \end{cases}                                      \\[6pt]
	x_{odjks} & = \begin{cases} 1 & \text{OD pair } (o,d) \text{ is assigned to pickup station } j \text{ and dropoff station } k \text{ in scenario } s \\ 0 & \text{otherwise} \end{cases}
\end{align*}

\paragraph{Objective.}
The objective minimizes total demand-weighted travel cost across all scenarios:
\begin{equation}
	\min_{x,y,z}
	\sum_{s \in S} \sum_{(o,d) \in \Omega} \sum_{(j,k)\in A_{od}}
	q_{ods}\,a_{odjks}\,x_{odjks}
	\label{eq:nominal-obj}
\end{equation}

where $a_{odjks} = d^{\mathrm{orig}}_{oj} + d^{\mathrm{dest}}_{dk} + \lambda\,c_{jk}$

\paragraph{Constraints.}
The constraints enforce the two-stage structure:
\begin{align}
	 & \sum_{(j,k)\in A_{od}} x_{odjks} = 1             &  & \forall (o,d), s \label{eq:assign} \\
	 & x_{odjks} \le z_{js}, \quad x_{odjks} \le z_{ks} &  & \forall (o,d), (j,k)\in A_{od}, s
	\label{eq:active}                                                                           \\
	 & z_{js} \le y_j                                   &  & \forall j, s \label{eq:build}      \\
	 & \sum_j y_j = l, \quad \sum_j z_{js} = k          &  & \forall s \label{eq:budget}
\end{align}
Constraint~\eqref{eq:assign} requires every OD pair to be assigned to exactly
one admissible station pair per scenario. Constraints~\eqref{eq:active} ensure
that assignments can only use active stations. Constraint~\eqref{eq:build}
requires that a station must be selected before it can be activated.
Constraint~\eqref{eq:budget} fixes the number of selected stations to $l$ and the
number of active stations per scenario to $k$.

\subsection{Robust Counterpart}

The nominal model assumes demand $q_{ods}$ is known exactly, but in practice it varies from day to day. To guard against this, we require the station selection to perform well under any demand realisation within a budget uncertainty set.

\subsubsection{Demand Uncertainty}

For each period $s$, demand lies in
\begin{equation}
	\cU_s = \Bigl\{
	q_s :
	\underline{q}_{ods} \le q_{ods} \le \overline{q}_{ods}\ \forall (o,d),\;
	\sum_{(o,d)\in\Omega} q_{ods} \le \overline{Q}_s, \; q_{ods} \in \mathbb{Z}_{+} \forall (o, d)
	\Bigr\}.
	\label{eq:uset}
\end{equation}

The per-OD bounds $[\underline{q}_{ods}, \overline{q}_{ods}]$ reflect the
observation that demand for a given OD pair fluctuates around a historical
range from day to day. There is also a aggregate limit defined by $\overline{Q}_s.$ which represents the idea that while demand can shift across OD pairs, the total number of trips within a time period is unlikely to increase beyond the population limit. Together, this uncertainty set is able to accurately represent worst case demand shifts over each scenario.

One concern is that since passenger demand is integral, this uncertainty set might be difficult to optimize over, but since the bounds $\underline{q}_{ods}, \overline{q}_{ods}, \overline{Q}_{s}$ should be integer values, the extreme points of the linear relaxation of the uncertainty set will also be integer. Thus, integrality can be dropped from the uncertainty set without changing the worst-case value.

\subsubsection{Robust Objective and Recourse}

Given fixed activation decisions $z_s$ and realized demand $q_s$, the optimal
assignment cost is captured by the recourse value:
\begin{equation}
	\Phi_s(z_s,q_s)
	:=
	\min_{x \in \{0,1\}}
	\sum_{(o,d)\in\Omega}
	q_{ods}
	\left(\sum_{(j,k)\in A_{od}}
	a_{odjks}\,x_{odjks}\right)
	\label{eq:recourse-value}
\end{equation}
subject to \eqref{eq:assign}--\eqref{eq:active}. Since this inner problem
decomposes by OD pair, letting
\[
	m_{ods}(z_s) = \min_{\substack{(j,k)\in A_{od} \\ z_{js}=1,\, z_{ks}=1}} a_{odjks}
\]
denote the cheapest admissible station pair where both pickup and dropoff are
active in scenario $s$, the recourse value simplifies to
\begin{equation}
	\Phi_s(z_s,q_s)
	=
	\sum_{(o,d)\in\Omega} q_{ods}\,m_{ods}(z_s).
	\label{eq:recourse-separable}
\end{equation}
The robust objective then minimizes the worst-case total cost across all
scenarios:
\begin{equation}
	\min_{y,z}
	\sum_{s \in S}
	\max_{q_s \in \cU_s}
	\Phi_s(z_s,q_s),
	\label{eq:robust-obj}
\end{equation}
subject to \eqref{eq:build}--\eqref{eq:budget}. As \eqref{eq:recourse-separable}
shows, $\Phi_s$ is linear in $q_s$ for fixed $z_s$, so the inner maximisation
over $q_s \in \cU_s$ is a linear program in $q_s$ and can be dualized to yield
a tractable reformulation.

\subsubsection{Derivation}

Introducing epigraph variables $\eta_s$, the robust problem becomes
\begin{equation}
	\min_{y,z,\eta} \sum_{s} \eta_s
	\label{eq:robust-epi}
\end{equation}
subject to
\begin{align*}
	 & \eta_s \ge \max_{q_s \in \cU_s} \Phi_s(z_s,q_s) &  & \forall s,   \\
	 & z_{js} \le y_j                                  &  & \forall j,s, \\
	 & \sum_j y_j = l,\quad \sum_j z_{js} = k          &  & \forall s,   \\
	 & y_j, z_{js} \in \{0,1\}.
\end{align*}

Substituting \eqref{eq:recourse-separable} and dualizing the inner problem gives dual
variables $\alpha_s \ge 0$, $\beta_{ods} \ge 0$, $\gamma_{ods} \ge 0$ such that
\begin{align}
	 & \eta_s \ge \overline{Q}_s\,\alpha_s
	+ \sum_{(o,d)} \overline{q}_{ods}\,\beta_{ods}
	- \sum_{(o,d)} \underline{q}_{ods}\,\gamma_{ods}
	 &                                                        & \forall s,
	\label{eq:robust-epi-dual}
	\\
	 & \alpha_s + \beta_{ods} - \gamma_{ods} \ge m_{ods}(z_s)
	 &                                                        & \forall (o,d),\, s.
	\label{eq:robust-con}
\end{align}

Expanding $m_{ods}(z_s)$ with the assignment variables $x$ and eliminating
$\eta_s$ yields the final MILP:
\begin{equation}
	\min_{x,y,z,\alpha,\beta,\gamma}
	\sum_s \overline{Q}_s\,\alpha_s
	+ \sum_s\sum_{(o,d)} \overline{q}_{ods}\,\beta_{ods}
	- \sum_s\sum_{(o,d)} \underline{q}_{ods}\,\gamma_{ods}
	\label{eq:robust-counterpart}
\end{equation}
subject to
\begin{align}
	 & \alpha_s + \beta_{ods} - \gamma_{ods} \ge \sum_{(j,k)\in A_{od}} a_{odjks}\,x_{odjks}
	 &                                                                                       & \forall (o,d),\, s,
	\label{eq:robust-con-final}
	\\
	 & \sum_{(j,k)\in A_{od}} x_{odjks} = 1
	 &                                                                                       & \forall (o,d),\, s,
	\label{eq:recourse-counterpart-assign-final}
	\\
	 & x_{odjks} \le z_{js},\quad x_{odjks} \le z_{ks}
	 &                                                                                       & \forall (o,d),\, (j,k)\in A_{od},\, s,
	\label{eq:recourse-counterpart-active-final}
	\\
	 & z_{js} \le y_j
	 &                                                                                       & \forall j,\, s,
	\label{eq:robust-build}
	\\
	 & \sum_j y_j = l,\quad \sum_j z_{js} = k
	 &                                                                                       & \forall s,
	\label{eq:robust-budget}
\end{align}
with $\alpha_s,\beta_{ods},\gamma_{ods}\ge 0$ and $x_{odjks} ,y_j, z_{js} \in \{0,1\}$.

\section{Experiments}

We use operational data from a microtransit system in Zhuzhou, China.
The full network has 84 stations, but for scalability we reduce to the 40 most
popular stations by OD request volume. Because these stations account for the
majority of demand, the reduction in order coverage is minimal.
\begin{table}[htbp]
	\centering
	\caption{Dataset variants for the April calibration and May backtest windows.}
	\label{tab:data-summary}
	\begin{tabular}{lrrr}
		\toprule
		Metric                & Full Zhuzhou & Sampled & Sparse \\
		\midrule
		Stations              & 84           & 40      & 40     \\
		Directed segments     & 6,972        & 1,560   & 1,560  \\
		April orders          & 2,741        & 2,560   & 2,560  \\
		April ODs with orders & 473          & 389     & 130    \\
		May orders            & 31,546       & 30,100  & 30,100 \\
		May ODs with orders   & 1,412        & 949     & 131    \\
		\bottomrule
	\end{tabular}
\end{table}

Table~\ref{tab:data-summary} describes the three variants. The sparse variant is constructed by randomly reassigning demand counts from the least popular OD pairs onto the most popular ones, concentrating all calibration demand onto the top 10\% of OD pairs (389 $\to$ 130). This mimics a data-poor setting where granular OD-level observations are unavailable and demand is only recorded at an aggregate level.

\subsection{Experimental Design}

We optimize station selection over data on April 1--3, 2025 and evaluate on May 1--31, 2025. We fix $l=30$ selected stations, a 300-second walking limit,
$\lambda=10$, and vary $k \in \{10,15,20\}$. For the robust model we vary the
calibration quantile $\tau \in \{0.90,0.95,0.99\}$.

\subsection{Uncertainty-Set Calibration}

Scenarios correspond to four time-of-day periods: morning (06--10h), afternoon
(10--15h), evening (15--20h), and night (20--24h).

\paragraph{Lower bounds $\underline{q}_{ods}$:}
$\underline{q}_{ods}=0$ for all OD pairs, since many pairs are absent on any
given day.

\paragraph{Upper bounds $\overline{q}_{ods}$:}
Every OD pair in the full network receives a strictly positive upper bound
$\overline{q}_{ods}$. For pairs observed at least once in the calibration
window, $\overline{q}_{ods}$ is the $\tau$-quantile of daily counts, treating
days with no demand for that pair as zero. For pairs with no observed orders,
$\overline{q}_{ods}$ is imputed via a rank-1 gravity model:
$\overline{q}_{ods} = \lceil O_s(o)\cdot D_s(d)/Q_s \rceil$, where $O_s(o)$
and $D_s(d)$ are the total origin and destination flows from observed pairs and
$Q_s$ is the total observed demand in period $s$. For stations that appear in
no orders at all, the gravity estimate falls back to a floor value equal to the
minimum positive $\overline{q}$ observed across all periods.

\paragraph{Aggregate cap $\overline{Q}_s$:}
$\overline{Q}_s$ is the $\tau$-quantile of total period demand across
calibration days.


\begin{table}[htbp]
	\centering
	\caption{Uncertainty-set diagnostics, 90th-percentile calibration.}
	\label{tab:uncertainty-summary}
	\small
	\begin{tabular}{lcccc}
		\toprule
		Period              & Sampled ODs & Sparse ODs & Sampled ratio & Sparse ratio \\
		\midrule
		Morning (06--10h)   & 94          & 80         & 17.42         & 17.46        \\
		Afternoon (10--15h) & 169         & 103        & 7.12          & 7.28         \\
		Evening (15--20h)   & 249         & 122        & 2.86          & 2.97         \\
		Night (20--24h)     & 120         & 87         & 13.16         & 13.31        \\
		\bottomrule
	\end{tabular}
\end{table}

Table~\ref{tab:uncertainty-summary} shows that both 40-station variants produce
large box-to-cap ratios, confirming that the aggregate cap $\overline{Q}_s$
substantially tightens the uncertainty set relative to a pure box. Without it,
the adversary could simultaneously push every OD pair to its upper bound, which
in the morning period corresponds to a scenario 17 times the observed total
demand, yielding an overly conservative solution.

\section{Results}

Both experiments report results on the full-OD dataset (949 OD pairs), covering
both the April calibration period (2,560 orders) and the May test period (30,100
orders). For the sparse experiment, this means evaluating on OD pairs not seen
during calibration: models were trained on 130 OD pairs while evaluation spans
all 949.

\subsection{Sampled-Network Results}

With dense training OD coverage, all models achieve zero uncovered orders
(\Cref{tab:three-model-summary}). The nominal model has the lowest mean cost for
every $k$; robust variants trade a small mean cost increase for lower variance.

\begin{table}[htbp]
	\centering
	\caption{%
		Sampled-network results, full-OD May orders.
		Bold = best in column within each $k$-block.%
	}
	\label{tab:three-model-summary}
	\small
	\resizebox{\textwidth}{!}{%
		\begin{tabular}{lllccccc}
			\toprule
			$k$ & $\tau$ & Model   & April mean cost  & May mean cost    & April std.       & May std.         & Uncovered  \\
			\midrule
			10  & --     & Nominal & \textbf{4266.09} & \textbf{4333.42} & 1129.72          & 1116.39          & \textbf{0} \\
			10  & 0.90   & Robust  & 4287.81          & 4359.76          & 1111.08          & 1102.67          & \textbf{0} \\
			10  & 0.95   & Robust  & 4302.24          & 4392.66          & \textbf{1108.89} & \textbf{1090.17} & \textbf{0} \\
			10  & 0.99   & Robust  & 4293.94          & 4369.71          & 1109.24          & 1098.48          & \textbf{0} \\
			\midrule
			15  & --     & Nominal & \textbf{3977.09} & \textbf{4049.94} & 1145.30          & 1119.97          & \textbf{0} \\
			15  & 0.90   & Robust  & 4018.31          & 4071.27          & 1106.48          & 1088.71          & \textbf{0} \\
			15  & 0.95   & Robust  & 4023.38          & 4080.47          & 1105.51          & 1085.45          & \textbf{0} \\
			15  & 0.99   & Robust  & 4030.38          & 4067.03          & \textbf{1096.25} & \textbf{1066.90} & \textbf{0} \\
			\midrule
			20  & --     & Nominal & \textbf{3924.95} & \textbf{3984.77} & 1168.67          & 1142.49          & \textbf{0} \\
			20  & 0.90   & Robust  & 3967.29          & 4003.44          & \textbf{1116.51} & \textbf{1082.44} & \textbf{0} \\
			20  & 0.95   & Robust  & 3948.96          & 3978.06          & 1118.75          & 1090.00          & \textbf{0} \\
			20  & 0.99   & Robust  & 3958.28          & 4005.45          & 1118.56          & 1091.22          & \textbf{0} \\
			\bottomrule
		\end{tabular}%
	}
\end{table}

\subsection{Sparse-Network Results}

The nominal model fits the sparse training data best but leaves 227 orders
uncovered at $k=10$ and 14 at $k=15$. Robust models achieve complete coverage
(\Cref{tab:sparse-summary}). At $k=20$ all models cover every order and costs
converge.

\begin{table}[htbp]
	\centering
	\caption{%
		Sparse-network results, full-OD May orders.
		Bold = best in column within each $k$-block.%
	}
	\label{tab:sparse-summary}
	\small
	\resizebox{\textwidth}{!}{%
		\begin{tabular}{lllccccc}
			\toprule
			$k$ & $q$  & Model   & April mean cost  & May mean cost    & April std.      & May std.         & Uncovered  \\
			\midrule
			10  & --   & Nominal & \textbf{4446.97} & \textbf{4305.57} & 851.46          & 1092.41          & 227        \\
			10  & 0.90 & Robust  & 4505.20          & 4363.35          & 845.97          & 1100.00          & \textbf{0} \\
			10  & 0.95 & Robust  & 4521.04          & 4392.62          & \textbf{836.33} & 1094.39          & \textbf{0} \\
			10  & 0.99 & Robust  & 4507.51          & 4368.66          & 840.04          & \textbf{1089.52} & \textbf{0} \\
			\midrule
			15  & --   & Nominal & \textbf{4182.48} & \textbf{4031.55} & 881.21          & 1082.65          & 14         \\
			15  & 0.90 & Robust  & 4234.00          & 4080.66          & 853.99          & 1059.06          & \textbf{0} \\
			15  & 0.95 & Robust  & 4235.88          & 4089.67          & \textbf{848.30} & \textbf{1058.01} & \textbf{0} \\
			15  & 0.99 & Robust  & 4236.58          & 4073.98          & 855.82          & 1067.54          & \textbf{0} \\
			\midrule
			20  & --   & Nominal & \textbf{4136.13} & \textbf{3972.53} & 901.73          & 1094.37          & \textbf{0} \\
			20  & 0.90 & Robust  & 4149.90          & 3985.01          & 886.97          & 1082.80          & \textbf{0} \\
			20  & 0.95 & Robust  & 4155.78          & 3998.38          & \textbf{883.00} & \textbf{1077.17} & \textbf{0} \\
			20  & 0.99 & Robust  & 4148.22          & 3984.41          & 887.77          & 1079.39          & \textbf{0} \\
			\bottomrule
		\end{tabular}%
	}
\end{table}

All three models select the same 30 stations (see \Cref{fig:built-k10} in the
appendix), so the divergence is entirely in which $k$ stations are activated per
period. \Cref{fig:activations-morning} shows the morning period at $k=10$: the
nominal model places all 10 activations in the high-demand core near the centre
of the network, while the robust ($q=0.95$) model distributes activations more
broadly, including peripheral stations that are rarely observed in the sparse
calibration data. This broader coverage is what eliminates the 227 uncovered
orders.

\begin{figure}[htbp]
	\centering
	\begin{subfigure}[b]{0.48\textwidth}
		\centering
		\includegraphics[width=0.9\textwidth]{../presentation/img/activation_nominal_period_1_k10.png}
		\caption{Nominal}
	\end{subfigure}
	\hfill
	\begin{subfigure}[b]{0.48\textwidth}
		\centering
		\includegraphics[width=0.9\textwidth]{../presentation/img/activation_robust_q095_period_1_k10.png}
		\caption{Robust ($q=0.95$)}
	\end{subfigure}
	\caption{Morning-period activations for $k=10$. Filled = active; hollow = selected but inactive. Nominal clusters all activations in the high-demand core; robust spreads them across the network.}
	\label{fig:activations-morning}
\end{figure}

However, for orders that \emph{are} covered, the cost distributions are nearly
identical. \Cref{fig:cdf-sparse-main} shows that CDFs for nominal and robust
largely overlap, with robust variants shifted only slightly to the right in the
upper tail. This reveals a gap between the expected and actual benefit of
robustness: the model was designed to hedge against worst-case costs, but in
practice the recourse function is simple enough that once a station is within
reach, the assignment cost is similar regardless of which model selected the stations.
The main benefit in the sparse setting is therefore service availability,
specifically ensuring that peripheral demand is not left uncovered, rather than cost
reduction for served orders.

\begin{figure}[htbp]
	\centering
	\begin{subfigure}[b]{0.48\textwidth}
		\includegraphics[width=\textwidth]{../experiments/2026-05-03_nominal_smoothed_robust_sparse/analysis/cdfs/cdf_k15_period_1.png}
		\caption{Morning (06--10h)}
	\end{subfigure}
	\hfill
	\begin{subfigure}[b]{0.48\textwidth}
		\includegraphics[width=\textwidth]{../experiments/2026-05-03_nominal_smoothed_robust_sparse/analysis/cdfs/cdf_k15_period_2.png}
		\caption{Afternoon (10--15h)}
	\end{subfigure}
	\caption{Cost CDFs for $k=15$, sparse experiment, May full-OD test set. Nominal and robust distributions largely coincide for covered orders.}
	\label{fig:cdf-sparse-main}
\end{figure}

\section{Conclusion}

A budgeted uncertainty set is a practical way to model uncertain OD demand
without extreme conservatism: per-OD bounds keep each pair within its historical
range while the cap $\overline{Q}_s$ prevents simultaneous surges. When data is
dense, robustness primarily reduces cost variance. When data is sparse, the
nominal plan concentrates activations in the observed high-demand core and leaves
peripheral orders uncovered, while the robust plan achieves full coverage at a
modest cost premium. Coverage can alternatively be enforced via explicit
constraints, which may be more efficient when variance reduction is not a goal.

Future work could incorporate richer recourse functions reflecting true
microtransit operations (vehicle capacity, rebalancing, multi-period dynamics),
and identify which demand regimes, such as low-density or spatially shifting
demand, most benefit from a robust plan.

\appendix
\section{Additional Figures}

\begin{figure}[htbp]
	\centering
	\includegraphics[width=\textwidth]{../experiments/2026-05-03_nominal_smoothed_robust_sparse/analysis/station_selections/station_selections_built_k10.png}
	\caption{Selected stations for $k=10$. Coloured dots = selected (30 of 40); gray = not selected. All three models select the same set of 30 stations.}
	\label{fig:built-k10}
\end{figure}

\begin{figure}[htbp]
	\centering
	\includegraphics[width=\textwidth]{../experiments/2026-05-03_nominal_smoothed_robust_sparse/analysis/station_selections/station_selections_activations_k10.png}
	\caption{Per-period activations for $k=10$ across all four time-of-day scenarios. Filled = active; hollow = selected but inactive; gray = not selected.}
	\label{fig:activations-all}
\end{figure}

\begin{figure}[htbp]
	\centering
	\begin{subfigure}[b]{0.48\textwidth}
		\includegraphics[width=\textwidth]{../experiments/2026-05-03_nominal_smoothed_robust_sparse/analysis/cdfs/cdf_k15_period_1.png}
		\caption{Morning (06--10h)}
	\end{subfigure}
	\hfill
	\begin{subfigure}[b]{0.48\textwidth}
		\includegraphics[width=\textwidth]{../experiments/2026-05-03_nominal_smoothed_robust_sparse/analysis/cdfs/cdf_k15_period_2.png}
		\caption{Afternoon (10--15h)}
	\end{subfigure}
	\vspace{0.5em}
	\begin{subfigure}[b]{0.48\textwidth}
		\includegraphics[width=\textwidth]{../experiments/2026-05-03_nominal_smoothed_robust_sparse/analysis/cdfs/cdf_k15_period_3.png}
		\caption{Evening (15--20h)}
	\end{subfigure}
	\hfill
	\begin{subfigure}[b]{0.48\textwidth}
		\includegraphics[width=\textwidth]{../experiments/2026-05-03_nominal_smoothed_robust_sparse/analysis/cdfs/cdf_k15_period_4.png}
		\caption{Night (20--24h)}
	\end{subfigure}
	\caption{Cost CDFs for $k=15$ across all four time-of-day scenarios, sparse experiment, May full-OD test set.}
	\label{fig:cdf-sparse-full}
\end{figure}

\end{document}
